\chapter{Blepharitis}

\section*{Definition and Overview}
Blepharitis is a common, chronic inflammation of the eyelids, specifically of the margins where the eyelashes grow. It usually affects both eyes and typically involves both upper and lower lids. The condition is uncomfortable and can cause crusting, redness, and a gritty sensation, but it does not usually damage eyesight. Two broad types are recognised: anterior blepharitis, which affects the front edge of the lid around the lashes, and posterior blepharitis, which relates to dysfunction of the meibomian (oil) glands along the lid margin. Blepharitis characteristically runs a relapsing course, flaring up and settling over many months or years.

\section*{Epidemiology}
Blepharitis is very common and is thought to affect up to a third or more of the population at some point in life. It is more prevalent in older adults and in people with oily skin, seborrhoeic dermatitis (dandruff), or rosacea. Posterior blepharitis, in the form of meibomian gland dysfunction, is one of the leading causes of dry eye syndrome. Although the condition is strongly associated with bacterial colonisation of the lids and skin disorders, it is not contagious and cannot be passed from person to person.

\section*{Aetiology and Pathophysiology}
Several interacting factors contribute to inflammation of the eyelid margin:

\begin{itemize}
    \item \textbf{Meibomian gland dysfunction:} the oil glands at the lid margins produce thickened, abnormal secretions that clog the glands and alter the tear film
    \item \textbf{Bacterial colonisation:} staphylococcal and other bacteria growing on the lid margin provoke inflammation
    \item \textbf{Skin disorders:} seborrhoeic dermatitis, rosacea, and eczema frequently involve the lids
    \item \textbf{Demodex mites:} tiny mites that inhabit the lash follicles of some people, releasing irritants
    \item \textbf{Allergy:} reactions to cosmetics, eye drops, or contact-lens solutions
    \item \textbf{Age-related changes} in gland structure and oil composition
\end{itemize}

These factors produce a chronic inflammatory state at the lid margin. The abnormal oil and inflammatory mediators destabilise the tear film, which is why dry eye so frequently accompanies blepharitis. Clogged glands can also give rise to styes (hordeola) and chalazia, and chronic inflammation may alter the lashes themselves.

\section*{Clinical Features}
Symptoms are bilateral, chronic, and worse in the morning:

\begin{itemize}
    \item Red, swollen, itchy, or burning eyelids
    \item A gritty, sandy, or foreign-body sensation in the eyes
    \item Crusting along the lid margins, with lashes that stick together, especially on waking
    \item Dandruff-like flakes or greasy scales at the base of the lashes
    \item Watery or dry eyes
    \item Sensitivity to light
    \item Eyelash problems, including loss, misdirection, or matting of lashes
    \item Recurrent styes or painless lid lumps (chalazia)
\end{itemize}

\section*{Differential Diagnosis}
The symptoms overlap with several other ocular conditions, and blepharitis frequently coexists with them:

\begin{itemize}
    \item Dry eye syndrome, which is commonly associated
    \item Conjunctivitis, whether viral, bacterial, or allergic
    \item Styes and chalazia
    \item Keratitis, or inflammation of the cornea
    \item Herpes simplex or herpes zoster infection of the eye
    \item Skin conditions affecting the lids, including eczema and psoriasis
    \item Rarely, an eyelid lesion or tumour that must not be overlooked
\end{itemize}

\section*{When to Seek Medical Care}
See a doctor or optometrist for a first episode of lid inflammation, for symptoms that persist despite two weeks of regular lid hygiene, or for any eye pain, increasing redness, or change in vision.

\textbf{Call 000 (or local emergency services) for:}
\begin{itemize}
    \item Sudden, severe eye pain or a rapid drop in vision
    \item A painful red eye with discharge, which may indicate a serious infection
    \item An eyelid that becomes very swollen, hot, and painful, suggesting spreading infection
    \item A foreign body in the eye that cannot be removed
\end{itemize}

\section*{Diagnosis and Investigations}
The diagnosis is clinical and is made by careful examination of the lids and eye surface.

\begin{itemize}
    \item \textbf{History:} duration of symptoms, associations with skin disease, use of cosmetics or contact lenses, and any previous treatments
    \item \textbf{Slit-lamp examination:} inspection of the lid margins, lash follicles, and eye surface under magnification
    \item \textbf{Tear-film assessment:} testing for accompanying dry eye
    \item \textbf{Gland expression:} gentle squeezing of the lid to assess the character of meibomian secretions
    \item \textbf{Swabs and cultures:} rarely needed, but useful for persistent or severe bacterial infection
    \item \textbf{Review of skin disease:} assessment of rosacea, seborrhoeic dermatitis, or eczema
    \item \textbf{Ophthalmology referral:} for severe, recurrent, or treatment-resistant cases
\end{itemize}

\section*{Management}
Blepharitis is managed with ongoing daily lid care rather than a single cure, and treatment aims to control symptoms and prevent flares.

\begin{itemize}
    \item \textbf{Warm compresses:} a clean cloth or compress held over the closed lids for several minutes to soften the oil
    \item \textbf{Lid massage:} gentle massage along the lid margins to express oil from blocked glands
    \item \textbf{Lid hygiene:} daily cleaning of the lid margins with a clean cloth, cotton bud, or proprietary lid-cleaner
    \item \textbf{Artificial tears:} lubricating drops to relieve accompanying dry eye
    \item \textbf{Topical treatments:} antibiotic or corticosteroid ointments on prescription for flare-ups
    \item \textbf{Managing skin disease:} treating rosacea or dandruff with medical and skincare advice
    \item \textbf{Professional therapies:} in-clinic meibomian gland expression, thermal pulsation, or intense pulsed light in resistant cases
    \item \textbf{Regular review:} to adjust treatment and monitor the eye surface
\end{itemize}

\section*{Complications}
\begin{itemize}
    \item Recurrent styes and chalazia
    \item Chronic dry eye, which can damage the corneal surface if severe
    \item Recurrent conjunctivitis
    \item Eyelash disorders, including loss and inward growth (trichiasis)
    \item Corneal complications from prolonged inflammation, including marginal keratitis or corneal ulcers in severe cases
    \item Cosmetic concerns, reduced quality of life, and time lost from work or study
\end{itemize}

\section*{Prognosis and Outcomes}
Blepharitis is a chronic condition without a permanent cure: most people control it well with daily lid hygiene, but it tends to flare intermittently and long-term maintenance is usually needed. With consistent care, complications such as styes, corneal problems, and severe dry eye are uncommon. Good control of any underlying skin condition such as rosacea or seborrhoeic dermatitis substantially improves the long-term outlook.

\section*{Prevention and Self-Care}
\begin{itemize}
    \item Make warm-compress and lid-hygiene routines part of daily life, even between flares
    \item Clean the lids gently and avoid rubbing the eyes
    \item Use fresh towels and wash hands before touching the eyes
    \item Avoid eye cosmetics during flare-ups, and replace old products regularly
    \item Remove all eye makeup and contact lenses before sleep
    \item Manage dandruff, rosacea, or eczema under medical advice
    \item Seek review at the first sign of a painful, red, or light-sensitive eye
\end{itemize}

\section*{Sources and Further Reading}
The following sources provide reliable, up-to-date information for patients and students of medicine.

\begin{itemize}
    \item NHS. \textit{Overview of blepharitis, its causes, and self-care.} https://www.nhs.uk/conditions/blepharitis/
    \item Mayo Clinic. \textit{Blepharitis -- symptoms, causes, and treatment.} https://www.mayoclinic.org/diseases-conditions/blepharitis/
    \item MSD Manual. \textit{Blepharitis -- professional edition.} https://www.msdmanuals.com/
    \item MedlinePlus. \textit{Blepharitis -- patient information.} https://medlineplus.gov/
    \item NICE. \textit{Clinical knowledge summaries on blepharitis.} https://cks.nice.org.uk/
\end{itemize}
